% Part: first-order-logic
% Chapter: beyond
% Section: other-logics

\documentclass[../../../include/open-logic-section]{subfiles}

\begin{document}

\olfileid{fol}{byd}{oth}

\olsection{منطقات أخرى}

كما لعلك أدركت الآن، ليس تصميم منطق جديد أمرًا صعبًا. فأنت أيضًا
تستطيع إنشاء تركيب خاص بك، وابتكار نسق استنباطي، وصوغ دلالة تلائمه.
وقد يلزمك شيء من البراعة إذا أردت أن يكون نسق ال!!{derivation} تامًا
بالنسبة إلى الدلالة، وقد يتطلب إقناع العالم بأن منطقك جدير حقًا
بالاهتمام بعض الجهد. لكنك تستطيع في المقابل قضاء ساعات ممتعة في
استكشاف خصائص منطقك الرياضية والحسابية.

شهدت العقود الأخيرة انفجارًا حقيقيًا في المنطقات الصورية. فالمنطق
الضبابي مصمم لنمذجة الاستدلال في الخصائص المبهمة، والمنطق الاحتمالي
مصمم لنمذجة الاستدلال في ظل عدم اليقين. ومنطقات الافتراض ومنطقات عدم
الرتابة مصممة لنمذجة صور قابلة للنقض من الاستدلال، أي استدلالات
«معقولة» يجوز نقضها لاحقًا عند ظهور معلومات جديدة. وثمة منطقات معرفية
مصممة لنمذجة الاستدلال في المعرفة؛ ومنطقات سببية مصممة لنمذجة
الاستدلال في العلاقات السببية؛ بل ومنطقات «معيارية» مصممة لنمذجة
الاستدلال في الواجبات الأخلاقية. وقد تجد هذه الكائنات مدروسة تحت عنوان
المنطق الرياضي، أو المنطق الفلسفي، أو الذكاء الاصطناعي، أو العلم
المعرفي، أو في موضع آخر، بحسب كون الدافع الأساس لإدخالها فلسفيًا أو
رياضيًا أو حسابيًا.

وتطول القائمة بلا نهاية، وتبدو الإمكانات غير محدودة. ولعلنا لا نحقق
أبدًا حلم لايبنتس برد العقل البشري كله إلى حساب---لكن ذلك لا يمنعنا
من المحاولة.

\end{document}
